%% kapa: exam=1 round=28 year=2017 period=1 form=A subject=civil_law subject_ko=민법 seq=7 number=7
%% source: Q-Net 감정평가사 기출문제 (cst00309) artlSeq=5199603 file=extracted/1008/5199603/2017년 제28회 감정평가사 1차시험 문제지/1교시 A형.pdf
%% conversion: kordoc 4.24.51 -> tools/gen1.py (auto); review=proofread-sonnet; flags=none
\begin{question}{7}
물건에 관한 설명으로 옳지 않은 것은? (다툼이 있으면 판례에 따름)
\begin{choices}
  \item 법률상 공시방법이 인정되지 않은 집합물이라도 특정성이 있으면 이를 양도담보의 목적으로 할 수 있다.
  \item 법정과실은 원칙적으로 수취할 권리의 존속기간 일수의 비율로 취득한다.
  \item 수목에 달려있는 미분리의 과실에 대해 명인방법을 갖추면 그 과실은 독립한 물건으로 거래의 목적으로 할 수 있다.
  \item 천연과실은 다른 특약이 있더라도 그 원물로부터 분리하는 때에 이를 수취할 권리자에게 속한다.
  \item 권원 없이 타인의 토지에서 경작한 농작물도 성숙하여 독립한 물건으로 인정되면 그 소유권은 명인방법을 갖출 필요 없이 경작자에게 있다.
\end{choices}
\end{question}
\answer{4}
